\section{Síntesis y ampliación}

Esta clase parte de tres trabajos de 2021 y propone un marco de modelado de audio que sigue vigente: primero se deciden la resolución temporal y la representation; después, el objetivo de predicción y la context hierarchy; por último, se evalúa con una metric adecuada a la tarea y se verifica si el modelo aprendió una estructura intermedia razonable. El Transformer es solo el sequence operator dentro de ese marco; lo que determina realmente las propiedades del sistema es cómo se conecta con la waveform, el spectrogram, el codebook, el pooling, la wavelet y el decoder.

\lecturefigure{slide-47-final-thoughts.jpg}{Conclusión de la clase: el Transformer es un avance importante, pero no debe considerarse el punto final de la evolución de las arquitecturas.}{Video oficial de Stanford Online 00:47:19--00:48:13.}

\teachervoice{Al final, el ponente dice que el Transformer parece «ahora puede resolverlo todo», pero espera que no sea el punto final, e invita al público a buscar la siguiente arquitectura de mayor impacto. La revisión del artículo cuatro años después y la evolución posterior de los modelos de audio confirman que se trata de una actitud de investigación y no de una declaración de victoria.}

\subsection{Los tres artículos condensados en un ciclo de ingeniería}

Los tres trabajos pueden unificarse en cinco pasos: primero, identificar la escala temporal que la tarea realmente requiere; elegir una representation de waveform, de espectro o de códigos; asignar un ancho de banda distinto al context local y al global; elegir una metric que permita poner a prueba una hipótesis en lugar de producir cifras vistosas; por último, verificar si la representación interna forma una estructura interpretable. Este ciclo es más reutilizable que «usar Transformer».

\begin{table}[H]
\centering
\small
\caption{Comparación unificada de las tres líneas de trabajo de esta clase.}
\begin{tabularx}{\textwidth}{L{0.17\textwidth}L{0.30\textwidth}L{0.18\textwidth}>{\raggedright\arraybackslash}X}
\toprule
Línea & Entrada y mecanismo de largo alcance & Métrica principal & Limitación más importante \\
\midrule
Generación de forma de onda cruda & Muestras discretas de 8 bits; attention local más condition comprimida & Predicción de muestras Top-5 & No es una métrica perceptual; el error de generación se acumula. \\
Aprendizaje de representaciones & VQ codes o vistas aumentadas; objetivo predictivo/contrastivo & Exactitud de sondeo lineal & El tokenizer y la augmentation determinan qué información se puede aprender. \\
Comprensión de audio & Patches de forma de onda aprendibles; jerarquía de pooling/wavelet & mAP multietiqueta & Datos y gama de líneas base limitados; no cubre el costo en streaming. \\
\bottomrule
\end{tabularx}
\end{table}

\subsection{Lista de errores frecuentes}

\begin{warningbox}{Seis errores de razonamiento que deben evitarse}
\begin{enumerate}
  \item Llamar «audio tokens» por igual a la raw waveform, al spectrogram y a los VQ token, sin precisar la token rate ni la pérdida de información;
  \item Usar la top-5 sample accuracy para inferir la calidad perceptual de música o voz;
  \item Usar una misma etiqueta $O(n^2)$ que ignora la estructura real de la local window, el downsampling y la hierarchy;
  \item Tomar la supervised/unsupervised linear-probe gap como una diferencia fija para todas las tareas posteriores;
  \item Afirmar que el modelo reproduce por completo la audición humana porque los learned filter se parecen a una sinusoid;
  \item Interpretar el título «Adieu» de 2021 como un rechazo de todas las arquitecturas convolucionales o híbridas posteriores.
\end{enumerate}
\end{warningbox}

\subsection{Conexión con cursos posteriores y sistemas modernos}

Este marco se extiende de forma natural a los neural audio codecs, las speech units, el audio-language pretraining, la generación por diffusion/flow, el streaming recognition y los modelos eficientes de secuencias largas. Al estudiar estos sistemas conviene plantear las mismas preguntas: cuál es el codec bitrate, si el tokenizer pierde phase/prosody, en qué escala interactúa el context, si el generador es paralelo y si la evaluación mide reconstruction, perception, semantic accuracy o human preference. Los nombres de los modelos nuevos no eliminan las restricciones de siempre.

\begin{importantbox}{Conclusión didáctica final}
Lo esencial del Audio Transformer no es que «attention también puede procesar sonido», sino lo siguiente: la representación convierte una señal física continua en una secuencia que se puede aprender; la estructura multiescala separa en niveles el detalle de alta frecuencia y la semántica de largo alcance; la función objetivo determina qué conserva la representación; y los límites de la métrica determinan qué podemos afirmar.
\end{importantbox}

\subsection{Lecturas adicionales}

\begin{itemize}
  \item \href{https://arxiv.org/abs/2106.16036}{Verma and Chafe, \emph{A Generative Model for Raw Audio Using Transformer Architectures}}: Transformer a nivel de muestra (sample-level) y context conditioning.
  \item \href{https://arxiv.org/abs/2010.11459}{Verma and Smith, \emph{A Framework for Generative and Contrastive Learning of Audio Representations}}: representation learning contrastive y de next-code.
  \item \href{https://arxiv.org/abs/2105.00335v1}{Verma and Berger, \emph{Audio Transformers} v1}: la raw-waveform understanding, el pooling, la wavelet y el learned front end que corresponden a la clase.
  \item \href{https://arxiv.org/abs/1609.03499}{van den Oord et al., \emph{WaveNet}}: línea base de dilated causal convolution.
  \item \href{https://arxiv.org/abs/1711.00937}{van den Oord et al., \emph{Neural Discrete Representation Learning}}: VQ-VAE y codebook learning.
  \item \href{https://arxiv.org/abs/2005.00341}{Dhariwal et al., \emph{Jukebox}}: hierarchical discrete music tokens.
  \item \href{https://arxiv.org/abs/2010.00475}{Fonseca et al., \emph{FSD50K}}: fuente del conjunto de datos de los experimentos de clasificación de la clase.
\end{itemize}

\end{document}
